\chapter{Implementation}\label{chap: application}

The theoretical concepts and methods described in \autoref{chap: Theory} and \autoref{chap: Methods} must be incorporated into a \acrshort{dns} simulation setup. 
Furthermore, the \acrshort{dns} data must be postprocessed, to extract the relevant quantities presented in \autoref{chap: Results}.
In the following, first, the simulation setup used in this theses is described in \autoref{sec: Simulation setup}, followed by a brief description of the postprocessing procedure in \autoref{sec: post-processing}.
Finally, in \autoref{sec: performed simulations} the performed simulations and chosen parameters are described.

\section{Simulation setup}\label{sec: Simulation setup}
For this study, the canonical setup of the pipe is used.
In the pipe, azimuthal periodicity is given by default meaning that a smooth pipe setup is fully defined using the pipe length and the Reynolds-number \citep{yaoDirectNumericalSimulations2023}.
Additionally, a cylindrical coordinate system is employed, with the three spatial directions $r, \theta$ and $z$.
The terminology wall-normal is used for the $r$-coordinate, azimuthal is the $\theta$ direction and $z$ is the streamwise direction.
The cartesian coordinates used by Neko can be transformed into the cylindrical coordinate system by the relations $r = \sqrt{x^2 + y^2}$ and $\theta = \arctan\!\left(\frac{y}{x}\right)$.

The dimensions of the rough domain are $L_r \times L_\theta \times L_z = R \times 2\pi R \times 8\pi R$, with the radius $R$ is set to $1 \hat{=}$ unity. 
The smooth validation domains are shortened to $L_z = 4\pi R$ to save computation time.
$R$ denotes the mean radius of the pipe in rough setups, and is therefore the distance from the virtual origin to the pipe center.

For the \glsentryfull{bcs}, the Dirichlet no-slip \acrshort{bc} at the pipe wall is used.
This states that the fluid at the wall  moves with the same velocity as the wall itself.
In the current setup all velocities are set to zero at the wall: $u_i = 0 \ \text{for}\ y = 0$.
This also implies that the wall is impermeable, since the wall normal velocity $v$ is zero.
Since the streamwise direction is statistically homogeneous in fully developed turbulent pipe flow \citep{moinDIRECTNUMERICALSIMULATION1998}, a periodic boundary condition is used for the inflow and outflow of the pipe.
This means that the outflow velocities are set equal to the inflow velocities: $u_{i_{l = 0}} = u_{i_{l=L_z}}$.
This implies that fluid leaving the domain at the outflow reenters the domain at the inflow.

It is important to use appropriate \glsentryfull{ics} to support a fast transition to turbulence.
This is favourable since it saves computational resources.
This is achieved by using a laminar velocity profile as a basis, which is overlaid by large- and small-scale disturbances.
The laminar velocity profile is defined as
\begin{equation}\label{eq: laminar velocity profile}
    u_r = 0 \quad u_\theta = 0 \quad u_z = 2(1-r^2)
\end{equation}
\begin{figure}
    \centering
    \includegraphics[width = 0.5\textwidth]{plots/Methods/Initial_condition-cropped.pdf}
    \caption[Visualisation of initial condition]{Visualisation of initial condition. Blue colours denote lower velocities, red denotes higher velocities.}
    \label{fig: initial condition}
\end{figure}
with $r$ denoting the radius ranging from $0$ to $1$.
The used profile represents the bulk velocity $U_b = 1$ in streamwise direction.
The perturbations are defined as a combination of multiple $\sin$ and $\cos$ waves of various magnitudes.
Large- and small-scale disturbances are intended to disrupt the laminar profile and support a fast transition to turbulence.
The initial velocity field is taken from the Neko example of the turbulent channel and converted to polar coordinates, and altered, such that the divergence-free criterion $\nabla u_i = 0$ is met, which is necessary for the conservation of mass \citep{popeTurbulentFlows2015}.
A visualisation of the initial velocity magnitude is depicted in \autoref{fig: initial condition}.
It can be seen, that the \acrshort{ic} is not physical, due to a presence of singularity at the pipe centre.
This is however unavoidable, since it was chosen to use the principle of sinusoidal disturbances in order to keep it in line with the \acrshort{ic} in the Neko channel example, but with a met divergence free criterion in polar coordinates.

The bulk flow rate $U_b$ is defined in the streamwise $z$ direction as unity ($U_b=1$) and is kept constant throughout the simulation.
To achieve the desired bulk velocity, the flow is driven by a pressure gradient $\Pi$.
To maintain a constant flow rate, the pressure gradient $\Pi$ is varying over time to adjust for changes in the flow conditions, e.g., when turbulence develops.
Using \autoref{eq: viscous scaling units} it is also possible to compute $Re_\tau$ from the pressure gradient $\Pi$  as described by~\cite{demaioDirectNumericalSimulation2023}
\begin{equation} \label{eq: REtau from forcing}
    Re_\tau = \sqrt{\frac{\Pi}{2}}Re_b \quad \text{since} \quad \tau_w = \frac{R}{2}\Pi 
\end{equation}

As mentioned in \autoref{chap: Introduction}, past numerical studies of turbulent flows over rough surfaces frequently used the \acrshort{ibm} method \citep{demaioDirectNumericalSimulation2023,busseReynoldsnumberDependenceNearwall2017,thakkarSurfaceCorrelationsHydrodynamic2017,frohnapfelFlowResistanceHeterogeneous2024}.
For the \acrshort{ibm}, a uniform mesh is constructed, which is slightly larger than the domain of interest.
Within the uniform mesh, the nodes that lie on the boundary between the desired shape and the flow domain are identified.
At each of the boundary nodes an additional forcing term is added, which enforces the \glsentryfull{bcs} \citep{verziccoImmersedBoundaryMethods2023}.
When comparing \acrshort{ibm} with body-fitted mesh topologies, both advantages and disadvantages emerge.
An advantage of the \acrshort{ibm} method is the significantly easier mesh topoglogy, which simplifies mesh generation.
Furthermore, because of the orthogonal mesh alignment, the computation effort per node is significantly reduced for the \acrshort{ibm} compared to a body-fitted mesh.
In contrast, the required number of grid points is higher for the \acrshort{ibm} method, with a difference with increasing Reynolds number.
The ratio of necessary grid-points for the \acrshort{ibm} to a body-fitted mesh can be estimated as $N_{IBM}/N_{BF}\approx Re^{0.9}$ \citep{verziccoImmersedBoundaryMethods2023}.
This results in a lower computation cost for the \acrshort{ibm} method for lower Reynolds numbers, due to the lower computation effort per node, but a higher computation cost for higher Reynolds numbers.
In this thesis a body fitted-mesh is chosen in order to provide data for wall roughness resolved using body-fitted meshes.

The simulation mesh must contain enough grid points to be able to correctly resolve the physics of turbulence.
Since the turbulent scales decrease near the wall and the viscous sublayer must to be resolved, it is necessary to refine the grid resolution near the wall.
Vice versa, the mesh can be coarsened towards the pipe center, without compromising the accuracy of the results.
\begin{figure}
    \centering
    \includegraphics[width = 0.5\textwidth]{Methods/mesh_19000Reb_crossSection.pdf}
    \caption[Visualisation of simulation mesh]{Exemplary visualisation of simulation mesh. Structured mesh topology in the outer part, unstructured mesh topology in the inner part.}
    \label{fig: simulation mesh}
\end{figure}
The simulation mesh is generated using Gmsh \citep{Gmsh}.
\autoref{fig: simulation mesh} shows the frontal view on one plane of an exemplary simulation mesh.
All the other meshes used in this work have a very similar topography.
It can be seen that the mesh consists of two sections.
The outer part is generated using a structured mesh topology, while the inner part uses an unstructured mesh.
This is done to avoid heavily distorted elements and singularity at the mesh center.
The size of the structured region is defined by an inner radius, which establishes the boundary between the structured and unstructured sections.
Increased resolution in the near-wall section is achieved through a geometric progression of the element size in the structured region of the mesh.
The number of elements in the azimuthal direction, $N_\theta$, also implicitly defines the number of elements in the unstructured part, since the size or number of mesh elements cannot be explicitly defined in the unstructured mesh topography.
All simulation meshes consist purely of quad elements.
To determine reasonable mesh resolutions, a mesh convergence study is performed in \autoref{subsec: grid refinement study}.
The resulting mesh parameters from this study are listed in \autoref{tab: parameters chosen meshes}.

For the time discretization, described in \autoref{subsec: time discretization}, the \acrshort{cfl}-number of 0.5 is chosen throughout this study, eventhough safe \acrshort{cfl} number range until $\approx 0.7$.
This choice is made, since the \acrshort{cfl}-numer varies throughout the simulation in a range of $\pm 0.1$.

\subsection{Applying rough surfaces onto the mesh}\label{subsec: Generating rough mesh}

The algorithm described in \autoref{subsec: Surface generation algorithm} yields the height values for a plane surface.
This surface has to be mapped onto the pipe in such a way, that the same generated surface can be mapped onto different meshes with varying amounts of grid points.
%This leads to the fact, that the straightforward method, i.e. applying the height values of the generated surface points directly onto the pipe mesh is not applied in this work.
The used approach relies on the fact that every surface topography can be decomposed into a multitude of $\sin$ and $\cos$ waves using a Fourier transform \citep{perez-rafolsGeneratingRandomlyRough2019}.
The decomposed values are the Fourier coefficients, which can be seen as a weight given for each mode.
By recombining the extracted modes, the surface can be recreated independently of the simulation mesh.
The Fourier coefficients are extracted by taking the real and imaginary part  of the Fourier transform of the generated rough surface, such that $A(m,n) = \operatorname{Re}\left(\hat{Z}_h\right)$ and $B(m,n) = \operatorname{Im}\left(\hat{Z}_h\right)$.

This decomposition and recreation is only exact, if an infinite number of modes is used.
%This is however only relevant for regular surfaces.
%Since the surfaces generated by the algorithm described in section \ref{subsec: Surface generation algorithm} are defined with a finite amount of modes (that is chosen by the user), this amount of modes is sufficient for the recreation of the surface.
With these parameters, the height function for the rough surface at each individual point can be defined as
\begin{equation}\label{eq: height function}
    k(\theta,z) = \sum_{m,n}^{M,N} A(m,n)\cos(phase)- B(m,n)\sin(phase)
\end{equation}
with
\begin{equation}\label{eq: phase equation}
    phase = 2\pi\left(\frac{uz}{L_z} + \frac{v\theta}{L_{\theta}}\right)
\end{equation}
with $A(m,n)$ and $B(m,n)$ representing the Fourier coefficients, $m$ and $n$ the Fourier mode and $z,\ \theta$ the streamwise and azimuthal coordinate of the respective point.
$M$ and $N$ denote the total number of Fourier modes that are taken into account.
The Fourier transform is performed using the Numpy python library, the result is ordered such that the first half of the array $A[1:M/2]$ contains the positive frequency terms, while the latter half $A[M/2+1:]$ contains the negative frequency terms, which follows the~\cite{NumpyFT}.
To get wave-numbers which are signed correctly, i.e.\ positive, $m$ and $n$ are adjusted accordingly, which gives the parameters $u$ and $v$.

The total number of Fourier modes $M$ and $N$ changes the shape of the surface significantly.
A higher number of Fourier modes also includes higher wave-numbers, which leads to surfaces with smaller roughness features. 
This however, resuts in a higher grid resolution requirement, since it shows that at least $12$ grid points per roughness feature are necessary, to sufficiently resolve the shape of the roughness feature \citep{busseDirectNumericalSimulation2015}.

The height function in \autoref{eq: height function} yields the recreation of the rough surface, which is applied on grid points at the wall.
However, a body fitted mesh is desired, which leads to the necessity, to create a scaling function for the grid points away from the wall.
The roughness function should be applied onto the mesh in the near wall region of the domain, while the inner region stays unaffected \citep{chanSystematicInvestigationRoughness2015}.
The used function for this purpose is
\begin{equation}\label{eq: roughness scaling function}
    r_{shift} = r + \frac{r - r_{crit}}{R - r_{crit}} * k(\theta,z) 
\end{equation}
with $r_{shift}$ representing the amount of necessary radial shift of the respective point, $r_{crit}$ the minimal radius, at which the grid points are shifted, $R$ the maximum radius of the smooth pipe mesh and $r$ the radial position of the individual point to be shifted.
This function scales the magnitude of shift applied to each grid point linearly between the outer pipe radius of the original smooth pipe and $r_{crit}$, from one to zero.
This method results in a body fitted mesh, where the boundaries of the mesh represent the roughness features.
Furthermore, the near wall mesh refinement is kept.

\begin{figure}
    \centering
    \includegraphics[width = 0.5\textwidth]{Methods/deformed_Mesh-cropped.pdf}
    \caption[Body fitted rough mesh]{Mesh deformed due to roughness}
    \label{fig: deformed mesh}
\end{figure}
The points are only shifted in radial direction.
Since the GLL points are distributed over the mesh before the roughness is applied, the GLL distribution is not exact after the rough deformation.
This effect is assumed to be negligable, since the shortening or stretching of the elements is of small magnitude.

\section{Postprocessing}\label{sec: post-processing}
\fxbmark{Diese section noch nicht}
In order to be able to work with the obtained data by the \acrshort{dns}, a post-processing of the data is necessary.
The \acrshort{dns} data is distributed on the \acrshort{gll} points and non uniformly distributed in space.
The post-processing is done, to map the data onto user defined points, since this simplifies all further processing steps.
All data is presented in a wall normal direction and after post-processing, the \acrshort{dns} data is mapped onto wall normal planes.

In this work, this is done using PySEMTools, which is an open source Python package, to post-process \acrshort{sem} data on hexahedral high-order elements \citep{perezPySEMToolsLibraryPostprocessing2026}.
Using PySEMTools it is possible to obtain fluid statistics from multiple instantaneous fields, which will be done in \autoref{chap: Results} for the energy spectra.
Additionally, Neko has the capability to produce statistics files, in which the data is averaged in time and, if chosen, in one or more spatial directions of choice during the runtime of the simulation.
These statistics files are used to obtain the relevant data for the first- and second order statistics plots.
In the present thesis, no spatial averaging is performed by Neko, so the statistics files contain 3D data that is only averaged in time.

The typical steps for post-processing are as follows:
\begin{itemize}
    \item In case multiple files exist that contain statistics data, average them to obtain a single time averaged field
    \item Interpolate the averaged data from the \acrshort{sem} mesh to user defined interpolation points
    \item Calculate the desired values, e.g.\ velocity profiles and turbulent budgets, based on the interpolated data
    \item Further average the 3D data in space, if desired
\end{itemize}
Special measures have to be taken to process rough data correctly.
The post-processing of rough data is described in the following section. 
The goal is to give a brief overview and explain the workflow in principle, thus some details are neglected.
\subsection{Postprocessing rough cases}\label{subsec: Postprocessing rough cases}
The post-processing of rough \acrshort{dns} data is not as straightforward as for smooth cases.
Since the virtual origin is not at the wall at most positions, the interpolation point-cloud  must account for that.


The approach used in the present thesis is to map the \acrshort{dns} domain, as well as the interpolation points into a computational domain.
In the computational domain, the simulation domain is mapped from the rough shape into a smooth pipe, while the interpolation points are mapped such that they are at the respective positions of concentric circles in the physical domain.
For the rough simulation domain, this is achieved by applying the inverse of the \autoref{eq: roughness scaling function} onto the rough mesh using the height functions \autoref{eq: height function} and \autoref{eq: phase equation}.
The inverse of \autoref{eq: roughness scaling function} is
\begin{equation}\label{eq: inverse height function}
     {r_{shift}}^{-1} = \frac{1.0}{\left(R + \frac{k(\theta, z)}{R - r_{crit}}\right)} * \left(r + \frac{r_{crit}}{R - r_{crit}} * k(\theta, z)\right)
\end{equation}

\begin{figure}
    \centering
    \begin{subfigure}{0.45\textwidth}
        \subcaption{}
        \includegraphics[width = \textwidth]{Methods/averaged_field.png}
        \label{subfig: averaged field}
    \end{subfigure}
    \hfill
    \begin{subfigure}{0.45\textwidth}
        \subcaption{}
        \includegraphics[width = \textwidth]{Methods/smooth_Averaged_field.png}
        \label{subfig: smooth averaged field}
    \end{subfigure}
    \plotverticalspacing
    \caption[Visualisation of the smoothing algorithm used in the postprocessing]{Visualization of the result of the smoothing algorithm. Color depicting the velocity magnitude of a time averaged field.~(a) Rough pipe domain, (b) Rough domain mapped to smooth pipe}
    \label{fig: comparison of smooth and normal averaged field}
\end{figure}
The results are shown from an exemplary case in \autoref{fig: comparison of smooth and normal averaged field}.
It can be observed, that the rough pipe is mapped back to the original radius of $R = 1$.
Yet, the data is scaled properly, such that the structure of the flow remains preserved.

To obtain correct statistics, the data should be interpolated on concentric circles in the physical domain.
\begin{figure}
    \centering
    \begin{subfigure}{0.325\textwidth}
        \subcaption{}
        \input{plots/Methods/interpolation_point_placement_1.pgf}
        \label{subfig: PP points at beginning}
    \end{subfigure}
    \hfill
    \begin{subfigure}{0.325\textwidth}
        \subcaption{}
        \input{plots/Methods/interpolation_point_placement_2.pgf}
        \label{subfig: PP points after mapping to smooth}
    \end{subfigure}
    \hfill
    \begin{subfigure}{0.325\textwidth}
        \subcaption{}
        \input{plots/Methods/interpolation_point_placement_3.pgf}
        \label{subfig: PP points after inverse maping}
    \end{subfigure}
    \plotverticalspacing
    \caption[Placement of data points and interpolation points while mapping in postprocessing]{Visualization of points mapping in post-processing. Dots \dnsdot~\ represent \acrshort{dns} data points, circles  \interpcircle~\ interpolation points, the  black line \wallline~\ the pipe wall and the grey area \domainsquare~\ the simulation domiain (a) Default point placement. (b) DNS domain mapped into the smooth computational domain, as depicted in \autoref{subfig: smooth averaged field}. (c) Interpolation points mapped by the inverse of the rough surface, such that they lie on the same respective concentric positions in the computational domain}
    \label{fig: Visualization of the PP points}
\end{figure} 
If the smooth averaged field would be interpolated onto concentric interpolation points, the generated data would be in a body fitted manner.
To achieve data extraction in a concentric way, the interpolation points have to be deformed as well. 
By applying \autoref{eq: inverse height function} onto the interpolation points as done for the rough pipe, the interpolation points are deformed in such a way, that they lie in the computational domain in the same positions, as they would be in the physical domain for a concentric data extraction. This principle is shown in \autoref{fig: Visualization of the PP points}.

\section{Performed simulations}\label{sec: performed simulations}

To improve the understanding of the behaviour of turbulence over rough surfaces, a multitude of simulations are performed in the present thesis.
The rough surfaces as well as the Reynolds numbers are altered, to gain a comprehensive insight into the roughness effects in turbulent pipe flows.
A total number of nine simulations are performed.
Three different surfaces are considered, and all three surfaces are studied using three different Reynolds numbers.
The chosen bulk Reynolds numbers $Re_b$ are $5300,\ 11700$ and $19000$.
Those were chosen, since they are the Reynolds numbers for the typical friction Reynolds numbers of past studies in smooth pipes of $Re_\tau = 180,\ 360$ and $550$ \citep{yaoDirectNumericalSimulations2023,elkhouryDirectNumericalSimulation2013,kimProgressPipeChannel2012}. 
The considered surfaces are all based on the study by~\cite{demaioDirectNumericalSimulation2023}.
It is chosen to use the typical bulk Reynolds-numbers rather than match the friction Reynolds-numbers, since $Re_\tau$ does not describe the turbulence behaviour as it does in smooth cases.
This is due to the fact, that it is not only dependent on the viscous drag, but also on the pressure drag in rough conditions.
This makes it also dependent on the surface shape and not only on the separation of scales, hence the viscosity, like it is in smooth cases.

\begin{table}
    \centering
    \caption[Surface parameters of simulation surfaces]{Surface parameters of simulation surfaces. $S1$ showing the greatest roughness height, $S2$ reduced roughness height with similar effective slope, $S3$ is $S1$ scaled down such that the roughness height of $S2$ is matched}
    \begin{tabular}{c c c c c c c}
        \toprule
                & $S_a$     & $S_q$     & $Sk$      & $S_{z_{max}}$ & $ES$      & $S_{z_{5 \times 5}}$  \\
        \midrule
        $S1$    & $0.031$   & $0.040$   & $-0.479$  & $0.306$       & $0.261$   & $0.228$               \\
        $S2$    & $0.021$   & $0.026$   & $-0.500$  & $0.230$       & $0.256$   & $0.171$               \\
        $S3$    & $0.024$   & $0.030$   & $-0.479$  & $0.229$       & $0.196$   & $0.171$               \\
        \bottomrule
    \end{tabular}
    \label{tab: surface parameters of production runs}
\end{table}
In \autoref{tab: surface parameters of production runs} the considered surface parameters for the three surfaces are listed.
The surface referred to as S1 resembles the grit-blasted surface of~\cite{demaioDirectNumericalSimulation2023} as a basis.
The $S1$ surface is not exactly the same, but an approximation.
The approximation is done, by creating a surface, that shows the same surface parameters, as the surface studied by~\cite{demaioDirectNumericalSimulation2023}.
It is created using the algorithm described in \autoref{subsec: Surface generation algorithm}.

Additionally, the surfaces $S2$ and $S3$ show a decreased surface height $k$ compared to $S1$.
$S2$ keeps the effective slope, despite having a decreased roughness height $k$.
\begin{figure}
    \centering
    \begin{subfigure}{0.8\textwidth}
        \centering
        \subcaption{}
        \subimport{plots/Methods/}{surface_height_distribution_piro.pgf}
        \label{subfig: surface height map piro}
    \end{subfigure}
    \\[-20pt]
    \begin{subfigure}{0.8\textwidth}
        \centering
        \subcaption{}
        \subimport{plots/Methods/}{surface_height_distribution_kDown.pgf}
        \label{subfig: surface height map kDown}
    \end{subfigure}
    \\[-20pt]
    \begin{subfigure}{0.8\textwidth}
        \centering
        \subcaption{}
        \subimport{plots/Methods/}{surface_height_distribution_scaled.pgf}
        \label{subfig: surface height map scaled}
    \end{subfigure}
    \begin{subfigure}{0.8\textwidth}
        \centering
        \hspace{0.5cm}
        \subimport{plots/Methods/}{heightmap_colorbar.pgf}
    \end{subfigure}
    \caption[Height maps of simulation surfaces]{Height maps of simulation surfaces. Red donting surface retractions, green denoting surface elevations.~(a) $S1$, (b) $S2$, (c) $S3$}
    \label{fig: surface height maps}
\end{figure}

$S3$ is a scaled-down version of $S1$, which results also in a decreased effective slope.
\autoref{fig: surface height maps} shows the height maps for all three surfaces.
Red and green shadings depict extreme values, while yellow shadings denote the mean height.
It can be seen, that $S1$ in \autoref{subfig: surface height map piro} shows the largest deviations from the mean height.
Furthermore, $S2$ in \autoref{subfig: surface height map kDown} exhibits more, but less pronounced roughness features compared to $S1$. 
$S3$ in \autoref{subfig: surface height map scaled} shows the same height map as $S1$ but with overall damped values.

It is chosen to use the same physical surfaces for all Reynolds numbers.
This leads to changing non-dimensionalised roughness heights $k^+$ between the Reynolds numbers.
This is done, since a fixed $h^+$ would necessitate the generation of new surface for every surface-Reynolds number combination.
However, due to the randomness and irregularity of the surfaces, this would lead to different physical shapes of the roughness features, even though the surface parameters remain the same. 
The similarity of flow conditions for surfaces with the same statistical surface parameters is not given, as discussed by~\cite{forooghiUniversalRoughnessCorrelation2017,chungPredictingDragRough2021}.
To mitigate this issue, the same surfaces were chosen for all Reynolds numbers.
This leads to different roughness heights in viscous length scales $k^+$ between the Reynolds numbers.

The cases with lowered roughness height $k$ show both the same roughness height.
It is chosen to be $\approx 75\%$ of the roughness height of $S1$.
This choice is made, since it is the lowest surface height with matched a $Es$, that can be generated using an appropriate number of Fourier modes that still enables the resolution requirement by~\cite{busseDirectNumericalSimulation2015}.
It can be seen, that all the surface parameters, that are comparable between the surfaces are within the margin of $0.05$, which is considered as sufficient.
The only exception is skewness, for which it proved difficult to find a good fit with the used algorithm.

